\documentclass[11pt]{article}
\usepackage[letterpaper,margin=1in]{geometry}
\usepackage{amsmath,amssymb,amsthm,mathtools,microtype}
\usepackage{xcolor,hyperref}
\hypersetup{colorlinks,linkcolor=blue!50!black,citecolor=blue!50!black,
  urlcolor=blue!50!black,pdfauthor={},
  pdftitle={The causal memory of gravitational radiation from compact area measurements}}
\newtheorem{theorem}{Theorem}[section]
\newtheorem{lemma}[theorem]{Lemma}
\newtheorem{proposition}[theorem]{Proposition}
\numberwithin{equation}{section}
\newcommand{\dd}{\mathrm d}
\newcommand{\RR}{\mathbb R}
\newcommand{\TT}{\mathbb T}
\DeclareMathOperator{\sech}{sech}
\DeclareMathOperator{\supp}{supp}
\DeclareMathOperator{\rank}{rank}
\DeclareMathOperator{\diag}{diag}
\setlength{\emergencystretch}{1em}
\title{The causal memory of gravitational radiation\\from compact area measurements}
\author{}
\date{}
\begin{document}
\maketitle
\begin{abstract}
We recover an outgoing gravitational load from a fixed compact spatial area profile on the retained side of a cut. The calibrated plane shear may have any smooth time-dependent potential, and the competing loads may also depend on time. All compact temporal drive and readout pairings on a window determine precisely the causal load on the window shortened at each end by the one-way travel time. A retarded Volterra transformation reduces the inverse problem to two convolution support tests. Exact Einstein--scalar backgrounds with winding rods, including anisotropic examples with nonzero Weyl curvature, realize the measured shear. Its outgoing load can have finite or infinite cross-window rank. For the inverse-square potential \(-\gamma/t^2\), that rank is exactly \(n\) when \(\gamma=n(n+1)\) and infinite for every other real \(\gamma\). The rank tests linear causal memory, rather than global spectral type or a Hamiltonian realization.
\end{abstract}

\section{Introduction}
Eliminating a propagating exterior replaces it by a causal boundary force law. This law need not be local in time. We call its dependence on the prior wall trace causal memory, distinct from asymptotic displacement memory, soft charges, or the corner data required by gravitational gauge symmetry \cite{P44}. For a laboratory confined to one side of the cut, two questions arise. Can it identify the exterior force law without sampling the wall? If so, can its data distinguish a finite linear state from an exterior with unbounded linear memory?

We give an exact answer for a calibrated, isolated gravitational shear. A fixed spatial profile, supported a positive distance from the cut, serves for both source and readout. Each source and readout is a difference of two spacelike area histories specified by material clocks and rods. The temporal tests range over all compact smooth gates in one finite window; no stationarity of the force law is assumed. The retained wave potential can vary arbitrarily in calibrated time. Causal transmutation and a support theorem identify the load on exactly the time square available after the two one-way flights.

The inverse theorem does not determine the memory rank in advance. The actual outgoing load can have rank one on a nonstationary material background, whereas another attained family has infinite rank. An inverse-square family gives a complete classification: the allowed finite ranks are all nonnegative integers, each at one specified coupling, with an exact realization of the nonlocal tail. The distinction can be read in every separated pair of windows within the recovered time square.

This combines the normal-force convention of \cite{P44} with a specific compact invariant area probe of the type studied in \cite{P47,Khavkine}. General volume-to-boundary transfer, time-dependent hyperbolic Dirichlet-to-Neumann inverse problems and material tensor stress already have established treatments \cite{BLOS,Kian,ENW}. Our result concerns one fixed positive-distance spatial profile, a sharp two-way temporal restriction, and an attained interacting Einstein--matter graviton channel. It assumes the retained calibration and does not reconstruct a general metric or all polarizations.

\section{A material area probe of the physical plane shear}
\label{sec:material}

Use Einstein gravity, with \(\kappa=8\pi G>0\), and a positive sigma model with potential,
\begin{equation}
 \mathcal L=\frac{R}{2\kappa}
 -\frac12(\nabla q)^2-\frac12 Z(q)\sum_{i=1}^3(\nabla\psi^i)^2
 -\mathcal V(q),\qquad Z>0 .
 \label{eq:L}
\end{equation}
The transverse directions \(x,y\) are periodic with coordinate lengths \(L_x,L_y\), while \(z\in\RR\). Take the metric and material fields to be
\begin{equation}
 g=c(t)^2(-\dd t^2+\dd z^2)+b(t)^2(\dd x^2+\dd y^2),\qquad
 q=q(t),\quad
 (\psi^1,\psi^2,\psi^3)=(\lambda x,\lambda y,\lambda_z z),
 \label{eq:metric}
\end{equation}
with \(q'(t)\ne0\). Thus each circle-valued material rod winds into a circle of the matching target period; no multivalued real scalar is used. By material labels we mean the fixed functions \(q_{\mathrm{bg}}^{-1}(q),\,\psi^1/\lambda,\,\psi^2/\lambda,\,\psi^3/\lambda_z\), with \(q_{\mathrm{bg}}^{-1}\) held fixed when varying the fields. We work in these labels. In the isotropic case \(c=b=a\), the parameter \(t\) is the calibrated conformal time, not FLRW proper time.

Torus averaging selects zero transverse momentum. Reflections of \(x\) and \(y\), each combined with reflection of its own material target angle, leave the background and action invariant. In this averaged sector the unique metric component odd under both is \(h_{xy}=h_{yx}=b^2\gamma_{xy}(t,z)\); all scalar fields and all other metric components transform with other characters. This makes the odd--odd linear equation closed, with no gauge or scalar modes discarded by hand. The potential has no mixed component in the trace-reversed equation \(R^\mu{}_\nu=\kappa(\partial^\mu q\,\partial_\nu q+Z\sum_i\partial^\mu\psi^i\,\partial_\nu\psi^i+\mathcal V\delta^\mu{}_\nu)\). The rods give the term \(-\kappa Z\lambda^2\gamma_{xy}/b^2\) there.

To see the geometric side explicitly, let \(Q=b^2(I+\gamma_{xy}\sigma_x)\) be the transverse metric block, where \(\sigma_x\) interchanges its two coordinates, and put \(B_\mu=Q^{-1}\partial_\mu Q\) for the two base directions. For a transverse block independent of \(x,y\) one has
\[
 R^a{}_b=-\tfrac12\nabla^\mu_{\rm base}(B_\mu)^a{}_b
         -\tfrac14\operatorname{tr}(B_\mu)(B^\mu)^a{}_b .
\]
Then \(B_t=2H I+\gamma_{xy}'\sigma_x\), \(B_z=(\partial_z\gamma_{xy})\sigma_x\), \(H=b'/b\), to first order. Since the two-dimensional base is conformally flat, this gives
\(\delta R^x{}_y=(\gamma_{xy}''+2H\gamma_{xy}'-\partial_z^2\gamma_{xy})/(2c^2)\).
Thus \(u=b\gamma_{xy}\) satisfies the actual calibrated wave equation and quadratic action
\begin{align}
 C(\partial_t^2-\partial_z^2+M(t))u&=j,
 & M(t)&=2\kappa Z(q)\lambda^2c^2/b^2-b''/b,\label{eq:wave}\\
 I^{(2)}[u]&=\frac C2\int (u_t^2-u_z^2-M(t)u^2)\,\dd t\,\dd z,
 & C&=L_xL_y/(2\kappa).\label{eq:act}
\end{align}
The kinetic measure before rescaling is \(b^2\,\dd t\,\dd z\); the Einstein tensor kinetic term counts the two entries \(xy,yx\) of the one symmetric mode. This fixes the absolute normalization \(C\) in the normal force. Changing the action by the time boundary term used in rescaling does not change the spatial wall force. No approximation by a fixed-background test scalar has entered the reduction.

The source \(j\) and the readings are relational area observables. At fixed labels let \(Q_{\rm sp}\) be the three-dimensional spatial metric and \(D=(\det Q_{\rm sp})Q_{\rm sp}^{-1}\) its cofactor. The areas \(a_\pm=\sqrt{v_\pm^{\mathsf T}D v_\pm}\) with coordinate normals \(v_\pm=e_x\pm e_y\) refer to spacelike marked planes. At the background \(a_\pm=\sqrt2\,bc\), and \(\delta D_{xy}=-c^2\delta(Q_{\rm sp})_{xy}\). Hence
\begin{equation}
 \delta(a_+-a_-)=-\sqrt2(c/b)\,\delta(Q_{\rm sp})_{xy}
     =-\sqrt2\,c\,u
 \label{eq:area}
\end{equation}
after normalized torus averaging. Off shell at first variation, the marker correction to the unmarked \(xy\) metric perturbation is \(b^2(\partial_x Z^y_{\rm mark}+\partial_y Z^x_{\rm mark})\), up to the sign convention for the displacement field. It integrates to zero on the torus, before any field equation. Thus multiplying the averaged area difference by the known fixed \(-1/(\sqrt2 c)\) defines the linear invariant probe. The variational drive with the same calibration has the same vertex. Any real product weight \(f(t)\varphi(r)\) is therefore a compact marked area source or reading for \(u\), with \(r=|z|\) on the chosen half-line. The cut is at the material label \(z=0\); the wave speed in \((t,r)\) is one. Measurements throughout concern the linearized sector with fixed calibration.

\section{The exact finite-window inverse}
\label{sec:recover}

Let \(I=(t_0,t_0+T)\) be a laboratory window contained in a slightly larger finite slab. On the retained side \(r=-z>0\), the calibrated potential \(M(t)\) is known and smooth and depends only on time. Write \(P_M=\partial_t^2-\partial_r^2+M(t)\). The constant \(C>0\), time, cut and normal coordinate are fixed. Choose one real nonzero \(\varphi\in C_c^\infty((0,\infty))\) and let \(d=\inf\supp\varphi>0\). The entire measured data for a load \(K\) are
\[
  \mathcal B_K(f,h)=\int h(t)\varphi(r)u_f(t,r)\,\dd t\,\dd r,
  \quad f,h\in C_c^\infty(I),\quad j(t,r)=f(t)\varphi(r).
\]
Here \(C\partial_r u(t,0)=K[u(\cdot,0)](t)\). Work with zero response before the drive and causal distributional kernels, with all indicated causal products defined. Specifically, if \(S_M\) is the retained Poisson map from trace to \(-\partial_r u|_0\), require that \(CS_M+K\) have a two-sided causal inverse \(D_K\) on the chosen zero-past domain. These assumptions hold for the outgoing extensions used below. Competitors need not be homogeneous, stationary, or described by a finite state; they are compared as effective scalar force laws for the same controlled sector, without an additional independently accessible drive port.

\begin{theorem}[One-profile recovery]\label{thm:recover}
If \(T>2d\), two such loads produce identical data on \(I\) if and only if their kernels agree on \(J\times J\), where
\[
                     J=(t_0+d,t_0+T-d).
\]
If \(T\le2d\), all such loads produce identical data on the open window. Contact terms of the load are included in the recoverable kernel. Time lengths refer to the calibrated unit-speed coordinate.
\end{theorem}

\subsection{Causal transmutation of the known retained propagation}
Fix a time \(a<t_0\) in the slab and require each history to be supported above some \(a'>a\). On this domain construct
\[
 (\mathcal T f)(t)=f(t)+\int_a^t k(t,s)f(s)\,\dd s,
\]
where \(k\) on \(a\le s\le t\) is the restriction of the Goursat problem
\begin{equation}
 k_{tt}-k_{ss}+M(t)k=0,\qquad
 k(t,t)=-\tfrac12\int_a^t M(\sigma)\,\dd\sigma,\qquad
 k(t,2a-t)=0 .
 \label{eq:goursat}
\end{equation}
Only the given \(M(t)\) for \(t\ge a\) is used. Extend the auxiliary kernel to \(2a-t\le s\le t\) for this construction, not the physical history. With \(u=(t+s)/2-a,\ v=(t-s)/2\) one solves
\[
 k(u,v)=g(u)-\int_0^u\int_0^v
      M(a+u'+v')k(u',v')\,\dd v'\,\dd u',
 \qquad g(u)=-\tfrac12\int_a^{a+u}M(\sigma)\,\dd\sigma .
\]
Here \(k(u,v)\) denotes the same kernel in characteristic coordinates. Picard terms are bounded by \(\|g\|_\infty(\|M\|_\infty uv)^j/(j!)^2\), and the integral equation gives smoothness by differentiation. The second-kind Volterra inverse of \(\mathcal T\) has an analogous smooth causal kernel: ordered integrals give factorial convergence on the slab. Both operators act on the indicated distributions by duality, with an irrelevant smooth cutoff below the lower support when a transpose is not compact near \(a\). Twice integrating by parts on smooth histories gives
\begin{equation}
  (\partial_t^2+M(t))\mathcal T=\mathcal T\partial_t^2 ,
  \qquad (\partial_t^2-\partial_r^2+M(t))\mathcal T
         =\mathcal T(\partial_t^2-\partial_r^2).
 \label{eq:transmute}
\end{equation}
Indeed the diagonal term in the difference is \(M(t)+2(k_t+k_s)(t,t)\), and the integral term is the interior expression in \eqref{eq:goursat}. The lower endpoint terms vanish by zero past; integration by parts extends the identity to the domain. This is an initial-value version of a classical transmutation construction \cite{KMT}. It is an exact causal map for retarded solutions and need not preserve a symplectic form or a quantum state.

In the massless reference a prescribed wall history \(g\) extends as \(g(t-r)\), so \(S_0=\partial_t\). Its readout against \(\varphi\) is \(W_0 g=\varphi*g\). For the zero-trace driven solution, the reflected retarded Green function similarly gives the inward wall force \(C\partial_r u_D|_0=\varphi*f\): normal differentiation gives the first arrival \(\delta(t-r')\) from a source at \(r'\). Since \(\mathcal T\) acts only in time, commutes with the spatial profile and boundary traces, and preserves causal solution uniqueness, both transferred operators for the chosen profile are
\begin{equation}
 S_M=\mathcal T\partial_t\mathcal T^{-1},\qquad
 W_M=\mathcal T W_0\mathcal T^{-1}.
 \label{eq:transfer}
\end{equation}
The two profile transfers coincide in this time-only class, by their common free reference. No general retarded reciprocity for time-varying spatial media is assumed. Conjugation enters the natural causal domain with a two-sided inverse and preserves the defined operator products by transport of that domain.

\subsection{Two support tests and the sharp window}
Split \(u_f\) into the known zero-trace response and the Poisson extension of the unknown wall trace. The zero-trace solution has inward force \(W_M f\); the Poisson extension contributes \(-CS_M\) times its trace, and the readout of that extension is \(W_M\) of its trace. Consequently, for the temporal source-to-readout operator,
\begin{equation}
       \mathcal R_K-\mathcal R_D=W_M D_K W_M,\qquad
       D_K=(CS_M+K)^{-1}.
 \label{eq:fac}
\end{equation}
Only the response \(\mathcal R_K\) is measured. The calibrated \(\mathcal R_D\) is known by solving the retained zero-trace problem.

Suppose two measurements agree. Conjugate \eqref{eq:fac} by \(\mathcal T^{-1}\) and set \(\widetilde K_i=\mathcal T^{-1}K_i\mathcal T\). Compact reference tests remain available even though \(\mathcal T\) need not preserve future compactness. For \(f\) compact in \(I\), test with a smooth cutoff of \(\mathcal T f\) which is equal to \(\mathcal T f\) up to any chosen \(T_1<t_0+T\) and dies inside \(I\). This is possible since \(\mathcal T f\) vanishes up to the starting time of \(f\). Causality permits removing the cutoff on that prefix and applying \(\mathcal T^{-1}\) to the output difference, which vanishes there and before the input. Letting the end of the prefix vary gives
\[
        W_0 \Delta W_0=0\text{ on }I^2,\qquad
   \Delta=(C\partial_t+\widetilde K_1)^{-1}-
                 (C\partial_t+\widetilde K_2)^{-1}.
\]
Write the following support argument after translating \(t_0\) to zero. For fixed test \(f\), the intermediate output \(\Delta W_0 f\) is supported above \(d+\inf\supp f>0\), and its convolution with \(\varphi\) vanishes until \(T\). The lower support identity of Titchmarsh--Lions \cite{Lions} gives \(\Delta W_0 f=0\) on \((0,T-d)\). Now fix \(h\in C_c^\infty((0,T-d))\). The tested remaining input distribution \(q(s)=\langle h(t),\Delta(t,s)\rangle\) is supported below \(\sup\supp h\) by causality. The reflected convolution \(\varphi(-\cdot)*q\) vanishes on \((0,T)\), and already has upper support below \(T\). The upper support identity yields \(\sup\supp q\le d\). These one-sided distributional support identities are local on the slab: one cuts off unused remote times before invoking the compact-support version. Varying product tests gives \(\Delta=0\) on \(J^2\).

The two-sided inverse identity \(L_1\Delta L_2=L_2-L_1\), \(L_i=C\partial_t+\widetilde K_i\), now gives equality of \(\widetilde K_i\) there. Causal products restricted to \(J^2\) use only intermediate times between source and observation, so depend only on their factors restricted there. This applies including boundary contact terms, and conjugating back gives equality of \(K_i\) on \(J^2\).

Conversely, equality of \(K_i\) on \(J^2\) gives equality of \(D_{K_i}\) there by the inverse identity in the opposite direction. The support of either transfer has delay at least \(d\), by \eqref{eq:transfer} and causal convolution. In fact whenever it lies in the slab, the first arrival of the single-pass map is exactly \(d\) after the input: the two causal inverses preserve the infimum of temporal support, and \(\varphi*\) adds \(d\) by the support identity, even when \(\varphi\) is flat at its support boundary. Hence both wall times in \eqref{eq:fac} necessarily lie strictly in \(J\); when \(T\le2d\), none are available. This proves Theorem~\ref{thm:recover}. It is exact uniqueness for all compact temporal gates, not a uniform stable inverse from noisy data.

\subsection{An outgoing local signature}
For a genuinely homogeneous temporal outgoing half-line, \(K_*=CS_M\), using the exterior \(M\) when it differs from that of the calibrated retained half-line. To extract its diagonal tail, let \(V=\mathcal T-I\). Integration by parts shows that \([V,\partial_t]\) is a smooth Volterra kernel with off-diagonal value \(-(k_t+k_s)(t,s)\), since the two coincidence terms cancel. Since \(S_M=\partial_t+[V,\partial_t]\mathcal T^{-1}\) and the product of two smooth Volterra kernels vanishes on the diagonal, 
\begin{equation}
 S_M(t,s)=\partial_t\delta(t-s)+{\boldsymbol1}_{s<t}H_M(t,s),
      \qquad H_M(t,t)=\tfrac12 M(t).
 \label{eq:diagTail}
\end{equation}
The diagonal value means the one-sided limit of the smooth tail after removing the contact term. Thus the recovered load reads out the exterior temporal potential in this class on the shorter window. It does not reconstruct a general unknown inhomogeneous exterior.

\section{Finite versus infinite causal memory on attained backgrounds}
\label{sec:rank}
For separated open subwindows \(J_-<J_+\), the rank of a load means the rank of its bilinear form on compact early drives and compact late tests. A system with \(m\) real linear internal coordinates and smooth time-varying local evolution has rank at most \(m\): the response to any early drive factors through the internal state at an intermediate time \cite{TVrank}. This argument includes locally available derivatives of the wall drive in the state equation. Differential contact terms contribute nothing between separated windows. A nonzero finite minor of the load certifies a lower bound on the necessary number of coordinates, and Theorem~\ref{thm:recover} makes this a test of the retained-side area data when both windows lie in \(J\).

\subsection{A constant positive calibrated mass}
For \(M(t)=m^2>0\), the actual outgoing half-line has \(S_{m^2}\) with Laplace multiplier \(\sqrt{p^2+m^2}\), and
\[
 S_{m^2}(t,s)=\partial_t\delta(t-s)
       +{\boldsymbol1}_{s<t}\,b_m(t-s),
       \qquad b_m(v)=mJ_1(mv)/v .
\]
The Laplace identity follows by differentiating the transform of \(mJ_1(mv)/v\) in \(p\), using \(\mathcal L[J_0(mv)]=(p^2+m^2)^{-1/2}\), and fixing the constant at \(+\infty\). Its analytic continuation as an entire function of \(v\) has
\begin{equation}
            b_m(v)=\frac1\pi\int_{-m}^m e^{i\omega v}
                              \sqrt{m^2-\omega^2}\,\dd\omega .
 \label{eq:besselFourier}
\end{equation}
At any distinct early times \(s_j\), a dependence among \(b_m(t-s_j)\) on the late interval continues by analyticity and would give \(\sum_j c_j e^{-i\omega s_j}=0\) for all \(\omega\) in an interval, by Fourier uniqueness and the positive weight. This exponential polynomial can vanish there only when all \(c_j=0\). Given finitely many independent translates, there are evaluation points in the late interval with a nonsingular matrix; otherwise all evaluation vectors lie in a common proper subspace. Compact smooth approximations to all the point gates preserve that minor. Therefore the outgoing rank is infinite on every separated time pair. Equation \eqref{eq:besselFourier} concerns the analytic tail, and is not the force spectral density.

\subsection{Exact nonstationary examples: the rank-one boundary}
Let \(B>0,\ 0\le\alpha\le B/4\), and let \(M_0\) be either \(0\) or \(B^2\). In \eqref{eq:metric} use \(Z=1,\ \lambda_z=\lambda,\ c=b=a\) and set
\[
 \kappa\lambda^2=\Lambda=\frac{B^2+M_0}{2},\quad
 w=\alpha\tanh(\alpha t),\quad
 a=e^{Bt}(B-w)/B,\quad \mathcal H=a'/a .
\]
The Darboux operator \(\mathcal A=\partial_t-w\) has \(w'+w^2=\alpha^2\), sends \(e^{Bt}\) to \(Ba\), and yields \(a''/a=B^2-2w'\). Define the homogeneous clock and its potential by the actual Einstein equations
\begin{equation}
 \kappa q'^2=2(\mathcal H^2-\mathcal H')-\Lambda,\qquad
 \kappa a^2\mathcal V(q)=\mathcal H'+2\mathcal H^2-\Lambda .
 \label{eq:FLRWsource}
\end{equation}
Indeed the Friedmann equation is \(3\mathcal H^2=\kappa(q'^2/2+3\lambda^2/2+a^2\mathcal V)\), and the spatial Ricci equation is the second equation in \eqref{eq:FLRWsource}. The rods are harmonic. With \(q'>0\), stress conservation gives the clock equation. The estimates
\[
 |\mathcal H-B|\le B/12,\qquad |\mathcal H'|\le 7B^2/144
\]
follow from \(\mathcal H=B-w'/(B-w)\), \(|w|\le\alpha\), \(0\le w'\le\alpha^2\), and \(|w''|\le2\alpha^3\). Since \(\Lambda\le B^2\), both \(q'^2\) and \(a^2\mathcal V\) have strictly positive lower bounds. The clock is globally increasing onto \(\RR\) and defines a smooth positive single-valued potential; the scalar curvature \(6(B^2-2w')/a^2\) is nonzero. These statements require no geodesic completeness or common potential for different parameters. At \(\alpha=0,\ M_0=B^2\) this is the coasting solution \(a=BT\) in proper time, with \(q=\kappa^{-1/2}\log T\) and \(\mathcal V=1/(\kappa T^2)\) after fixing a clock offset. For \(\alpha>0\) the expansion is noncoasting, and the calibrated shear potential from \eqref{eq:wave} is
\begin{equation}
                          M(t)=M_0+2w'(t).
 \label{eq:Dmass}
\end{equation}

On zero-past distributions the differential \(\mathcal A\) has the two-sided causal inverse
\[
 (\mathcal A^{-1}f)(t)=\cosh(\alpha t)
                 \int_{-\infty}^t \sech(\alpha s)f(s)\,\dd s.
\]
The homogeneous integration constant is ruled out by the zero past. The Darboux identity \(P_{M_0+2w'}\mathcal A=\mathcal A P_{M_0}\) follows from \(-w''-2ww'=0\), and commutation with spatial differentiation gives the exact outgoing load \(CS_M=C\mathcal A S_{M_0}\mathcal A^{-1}\). Applying this causal map does not identify quantum states or act as a pointwise change of units. This use of the classical Darboux factorization \cite{PT} evaluates the physical force law, while Section~\ref{sec:recover} does not require such a special potential.

For \(M_0=B^2\), the map \(f\mapsto S_{M_0}f|_{J_+}\) has infinite-dimensional image for early compact drives. The new early drives \(\mathcal A f\) are again compact there and produce \(\mathcal A S_{M_0}f\) on \(J_+\). The kernel of \(\mathcal A\) on a connected late interval has dimension one, so the new image is infinite dimensional. The actual kernel is nonstationary for \(\alpha>0\) already by its tail diagonal \eqref{eq:diagTail}.

For \(M_0=0,\ \alpha>0\), use \(S_{M_0}=\partial_t\) and \([\mathcal A,\partial_t]=w'\). One obtains instead
\begin{equation}
 CS_M(t,s)=C\partial_t\delta(t-s)
 +C\alpha^2\,{\boldsymbol1}_{s<t}\sech(\alpha t)\sech(\alpha s).
 \label{eq:rankOne}
\end{equation}
It has rank exactly one on every separated pair; a state with
\(X'=\sech(\alpha t)\,g(t)\) and memory output \(C\alpha^2\sech(\alpha t)X(t)\) realizes its full smooth tail. Its positive, genuinely time-dependent wave potential and unbounded radiative half-line are compatible with finite linear causal memory in this channel. A finite tail plus derivative contact does not exhibit a finite closed lossless model of the entire exterior.

\subsection{Anisotropic examples with nonzero Weyl curvature}
The two Darboux potentials are also attained without conformal flatness on any prescribed finite time slab. Use \(Z=1\) in \eqref{eq:metric}, \(H=b'/b,\ \rho=c/b,\ \ell=\rho'/\rho\), and prescribe the target \(M(t)\) in \eqref{eq:Dmass}. Integrate
\begin{align}
  H'&=2\kappa\lambda^2\rho^2-M-H^2,& b'&=Hb,\notag\\
  \rho'&=\rho\ell,&
  \ell'&=-2H\ell-\kappa(\lambda^2\rho^2-\lambda_z^2),
  \label{eq:bianchi}
\end{align}
and set
\begin{align}
 \kappa c^2\mathcal V(q)&=H'+2H^2-\kappa\lambda^2\rho^2,\notag\\
 \kappa q'^2&=2(H^2-H')+4H\ell-\kappa\lambda_z^2.
 \label{eq:bianchimatter}
\end{align}
The exact spatial Riccis in the orthonormal frame are \(c^2R_{\hat x\hat x}=H'+2H^2\) and \(c^2R_{\hat z\hat z}=(H+\ell)'+2H(H+\ell)\). Their difference is the anisotropy equation in \eqref{eq:bianchi}. The Hamiltonian constraint is \(c^2G_{\hat0\hat0}=3H^2+2H\ell\); substituting the spatial Einstein equation gives the second line of \eqref{eq:bianchimatter}. The first line gives the matter potential, the harmonic rods remain exact, and Bianchi supplies the homogeneous clock equation. Equation \eqref{eq:bianchi} for \(H'\) fixes precisely the desired shear mass. These equations therefore prescribe a self-consistent gravity--matter solution, not an externally imposed anisotropic stress.

Choose \(\lambda_z=\lambda\), and at one interior time match \(b,H\) to the FLRW solution with \(\rho=1\), but take a nonzero sufficiently small \(\ell\). The FLRW reference has strict lower bounds for clock and potential on a slightly larger compact interval. ODE continuous dependence preserves the bounds and the positive scales for the perturbed exact solution on the slab. The resulting potential can be positively extended beyond the realized clock range; no common matter potential across the family or global anisotropic stability is asserted. In the curvature convention of \eqref{eq:FLRWsource},
\[
 C_{\hat0\hat x\hat0\hat x}-C_{\hat0\hat z\hat0\hat z}
        =\ell'/(2c^2).
\]
Indeed the Riemann difference times \(c^2\) is \(\ell'+H\ell\), and the Ricci subtraction is \(-(\ell'+2H\ell)/2\). At the chosen time \(\ell'=-2H\ell\ne0\). This is an intrinsically nonconformally-flat example, with the same isolated shear action, force calibration and compact area experiment from Section~\ref{sec:material}. The existence on finite slabs is sufficient for the causal test; past continuation of this perturbed background is not used.

\subsection{Generic infinite rank in the time-only class}
The exceptional finite-rank examples do not persist on any open set of smooth temporal potentials. Fix a separated pair and a positive constant reference. For each integer \(N\), pick compact gates with a nonzero \(N\)-minor at the reference. The same minor for \(M_z=(1-z)M+zM_{\rm ref}\) is an entire function of \(z\): the Goursat and Volterra series on the fixed slab, and the first kernel derivatives in \eqref{eq:diagTail}, converge locally uniformly in this parameter. In particular \((k_t+k_s)=\partial_u k\) in the characteristic coordinates of the Goursat problem, so uniform control of \(M\) suffices for continuity of a fixed minor. The entire function is nonzero at \(z=1\), and hence its real zeros are isolated. The nonvanishing condition is open and dense. Intersecting these conditions for each \(N\) and for a countable rational basis of separated interval pairs gives a residual set of real smooth \(M\) with infinite rank on every such pair. This genericity is a statement in the calibrated temporal operator class, not across all Einstein solutions in a single fixed matter theory. A finite witness is stable against small changes of that finite load minor, which is different from a uniform noisy raw-data inverse.

A neighborhood of such temporal potentials is also physically attained if the clock potential is allowed to vary. In the isotropic model fix \(Z=1,\ 2\kappa\lambda^2=B^2\) and solve \(\mathcal H'=B^2-M-\mathcal H^2\), starting at \(\mathcal H=B\) at an interior time, with \(a'/a=\mathcal H\). At \(M=0\) the positive clock and potential margins in \eqref{eq:FLRWsource} are strict. For any sufficiently small smooth temporal \(M\) on a fixed slightly enlarged slab, ODE continuity preserves these margins. Equation \eqref{eq:FLRWsource} then gives an exact background and a positive potential on the realized clock range. The genericity statement restricts to this open attainable set of potentials; the action is not held fixed across it.

\section{Exact memory ranks for an inverse-square shear}
\label{sec:ladder}

\begin{theorem}[Inverse-square classification]\label{thm:ladder}
Let \(M(t)=-\gamma/t^2\) for \(t>0\) and real \(\gamma\), and take the actual outgoing half-line with initially zero histories supported away from \(0\). On every separated pair of nonempty positive-time windows its calibrated load has rank exactly \(n\) if \(\gamma=n(n+1)\) for \(n\in\{0,1,\ldots\}\). At every other real \(\gamma\) it has infinite rank. At the finite-rank values there is an explicit minimal \(n\)-coordinate linear realization after retaining the local time derivative contact. Every \(\gamma\) in this statement is attained by an exact positive-kinetic four-scalar Einstein background with the marked area probe.
\end{theorem}

\subsection{The retarded layer and its logarithmic-time transfer}
For \(\nu(\nu+1)=\gamma\), the whole-line retarded solution from \((s,0)\) has kernel
\begin{equation}
 G(t,r;s)=\tfrac12{\boldsymbol1}_{t-s\ge|r|}
          P_\nu\!\left(\frac{t^2+s^2-r^2}{2ts}\right).
 \label{eq:LegGreen}
\end{equation}
The Legendre branch is analytic at \(1\), with \(P_\nu(1)=1\). This is the classical Riemann function for the inverse-square potential \cite{Legendre}. To check the normalization, write \(\xi=(t^2+s^2-r^2)/(2ts)\). It obeys \(\xi_t^2-\xi_r^2=(\xi^2-1)/t^2\) and \(\xi_{tt}-\xi_{rr}=2\xi/t^2\); hence the Legendre equation gives the interior wave equation. On each characteristic edge \(\xi=1\). The tangential derivative terms from the retarded step functions vanish there, and the surviving corner delta equals that of the free \(1+1\) retarded kernel. The two roots \(\nu\) give the same real analytic normalized function by its equation and regular branch.

Take the single layer \(u(t,r)=\int G(t,r;s)f(s)\,\dd s\) for \(r>0\). The normal derivative at \(0^+\) is \(-f(t)/2\): differentiating the upper time \(t-r\) gives that value, while the derivative of the smooth interior factor tends to zero on positive-time compact slabs. With wall value \(g\) and outgoing \(S_\gamma g=-\partial_r u|_0\) this gives
\[
          g(t)=\int_0^t P_\nu(\cosh(\log(t/s)))\,S_\gamma g(s)\,\dd s .
\]
The time integral is effectively over \(s>0\) bounded away from zero for each zero-past input. Differentiating this Volterra equation with unit diagonal gives the unique causal density, so the layer constructs the outgoing solution for the prescribed trace. Only finite-time uniqueness is required.

Set \(\tau=\log t,\ \widetilde g(\tau)=g(e^\tau),\ Y(\tau)=e^\tau S_\gamma g(e^\tau)\), and \(L(v)=P_\nu(\cosh v)\). Then \(\widetilde g=({\boldsymbol1}_{v\ge0}L)*Y\), with Laplace transfer \(1/\widehat L(p)\) from \(\widetilde g\) to \(Y\).
The analytic tail \(h_\gamma\) of this inverse is defined by
\begin{equation}
 \mathcal L^{-1}\bigl(1/\widehat L\bigr)
       =\delta'+{\boldsymbol1}_{v>0}h_\gamma(v),\qquad
 S_\gamma g(t)=g'(t)+\int_0^t\frac{h_\gamma(\log(t/s))}{ts}\,g(s)\,\dd s .
 \label{eq:physicaltail}
\end{equation}
Here \(\widehat L(p)=\int_0^\infty e^{-pv}L(v)\,\dd v\). To justify the analyticity and transforms, the derivative of the causal kernel is \(\delta+{\boldsymbol1}_{v>0}L'\), since \(L(0)=1\). Its Volterra inverse is \(\delta+{\boldsymbol1}_{v>0}R\), where \(R(0)=-L'(0)=0\). Its Neumann series, formed with the holomorphic continuation of \(L'\) on a small convex neighborhood of any real compact delay interval \([0,L_0]\), converges uniformly on compact subneighborhoods: the \(j\)-fold ordered integrals have bound \(C^j |v|^{j-1}/(j-1)!\). Thus \(h_\gamma=R'\) is analytic on positive delays with no additional delta. \(L,L',L''\) grow at most exponentially by \(L''+\coth v\,L'-\gamma L=0\) where \(v\ge1\); the Volterra series and its differentiated equation likewise bound \(R,R'\) exponentially. All Laplace identities hold on a common right half-plane. In physical time the derivative contact is \(g'(t)\), and the factors \(1/(ts)\) follow from changing both the output and integration clock. The diagonal value in \eqref{eq:physicaltail} is \(-\gamma/(2t^2)\), also agreeing with \eqref{eq:diagTail}.

\subsection{Integer couplings and the converse}
For \(\gamma=n(n+1)\), the generating function
\[
 (1-2u\cosh v+u^2)^{-1/2}
       =(1-u e^v)^{-1/2}(1-u e^{-v})^{-1/2}
\]
gives \(L(v)=\sum_{j=0}^n c_j e^{(n-2j)v}\) with \(c_j>0\), \(\sum_j c_j=1\), and \(c_j=c_{n-j}\). The \(n+1\) simple poles of \(\widehat L\) have positive residues. Between every two consecutive poles the transform is strictly decreasing from \(+\infty\) to \(-\infty\); its degree-\(n\) numerator has exactly \(n\) simple interlacing zeros. Their positions are explicit. The Legendre ODE, multiplied by \(2\sinh v\) and Laplace transformed with \(L(0)=1,L'(0)=0\), gives
\[
 (p^2-p-\gamma)\widehat L(p-1)=(p^2+p-\gamma)\widehat L(p+1).
\]
For integer \(n\), the rational function
\begin{equation}
 \widehat L(p)=
   \frac{\displaystyle\prod_{j=0}^{n-1}[p-(n-1-2j)]}
        {\displaystyle\prod_{j=0}^n[p-(n-2j)]}
 \label{eq:rational}
\end{equation}
satisfies this identity by cancellation and behaves as \(1/p\). Any two such rational solutions have a rational periodic quotient (period two), hence that quotient is constant, fixed by the normalization. Write \(\mu_j=n-1-2j\) and \(r_j=1/\widehat L'(\mu_j)<0\) for \(0\le j<n\). Then the inverse has Laplace multiplier \(p+\sum_j r_j/(p-\mu_j)\), with no constant contact. The real load tail is
\begin{equation}
       CS_\gamma(t,s)\big|_{t>s}=C\sum_{j=0}^{n-1}r_j\,t^{\mu_j-1}s^{-\mu_j-1}.
 \label{eq:laddertail}
\end{equation}
These powers are independent on each of the separated intervals because distinct exponentials are independent in logarithmic time. Evaluation minors can be realized with compact smooth gates as in Section~\ref{sec:rank}. The rank is exactly \(n\), including the zero-tail case \(n=0\). The minimal realization in logarithmic time is \(\dot X_j=\mu_j X_j+\widetilde g,\ Y=\partial_\tau\widetilde g+\sum_j r_j X_j\), with zero past. The \(\mu_j\) are transfer exponents, not graviton energy poles.

For the converse, finite rank of the kernel on any separated rectangle is unchanged by the invertible clocks and weights in \eqref{eq:physicaltail}. The translates \(h_\gamma(\tau-\sigma)\) with \(\sigma\) in the early interval must span a finite-dimensional space on the late interval. Otherwise evaluation minors, and then nearby smooth gate minors, would witness unbounded rank. Differentiating the translations at an interior \(\sigma_0\) stays in this finite-dimensional space, by closure under limits on a compact late subinterval. Hence a constant-coefficient differential equation holds for \(h_\gamma\) on one open delay interval and, by real analyticity, on all positive delays. Thus \(h_\gamma\) is an exponential polynomial. Its rational Laplace transform makes \(1/\widehat L\), and hence \(\widehat L\), rational. The latter is strictly proper because \(L(0)=1\). Partial fractions and Laplace uniqueness then make \(L\) an exponential polynomial as well.

Suppose \(L=\sum_\lambda a_\lambda(v)e^{\lambda v}\) with nonzero polynomials \(a_\lambda\) and distinct complex exponents. The ODE becomes
\[
 e^v(D^2+D-\gamma)L-e^{-v}(D^2-D-\gamma)L=0 .
\]
Group exponents by differences in \(2\mathbb Z\), the only possible cross-term cancellations after the shifts. In a nonempty class, the top exponent satisfies \(\lambda_{\max}^2+\lambda_{\max}-\gamma=0\) by its unmatched top term and the leading coefficient of its polynomial. Similarly the bottom satisfies \(\lambda_{\min}^2-\lambda_{\min}-\gamma=0\). Their difference must be an even nonnegative integer. The respective roots are \(\{\nu,-\nu-1\}\) and \(\{-\nu,\nu+1\}\), giving differences \(2\nu,-1,-1,-2\nu-2\). Therefore \(\nu=n\) or \(-n-1\) for an integer \(n\ge0\); the repeated-root and complex-root cases require no exceptions. This proves the analytic rank classification.

\subsection{Attainment by matter}
In \eqref{eq:L} take FLRW proper time \(T_{\rm p}>0\), \(a=T_{\rm p}^{\,q_0}\) with \(1/2<q_0<1\), and unit winding slopes \(\lambda=\lambda_z=1\). Write \(M_{\rm Pl}^2=1/\kappa\). For a constant \(\sigma>0\) set
\begin{align*}
 q&=k\log T_{\rm p},& k^2&=2M_{\rm Pl}^2q_0-\sigma,\\
 Z(q(T_{\rm p}))&=\sigma\,a^2/T_{\rm p}^2,&
 \mathcal V(q(T_{\rm p}))&=[M_{\rm Pl}^2(3q_0^2-q_0)-\sigma]/T_{\rm p}^2 .
\end{align*}
A constant reference unit in the logarithm or scale factor is absorbed in these definitions. Choose \(0<\sigma<M_{\rm Pl}^2 q_0(3q_0-1)\) and \(k>0\). This bound gives positive potential and a positive kinetic target metric and is stricter than \(k^2>0\). The energy density is \(3M_{\rm Pl}^2q_0^2/T_{\rm p}^2\) and \(\rho+p=2M_{\rm Pl}^2q_0/T_{\rm p}^2\), exactly solving the Friedmann and acceleration equations. The weighted rods \(\nabla_\mu(Z\nabla^\mu\psi^i)=0\) are harmonic by spatial homogeneity. Their conservation and nonzero clock derivative give the clock equation, or one substitutes into
\(\ddot q+3(q_0/T_{\rm p})\dot q+3Z_{,q}/(2a^2)+\mathcal V_{,q}=0\), with proper-time derivatives here. Scalar curvature is \(6q_0(2q_0-1)/T_{\rm p}^2>0\).

The calibrated conformal time is \(t=T_{\rm p}^{\,1-q_0}/(1-q_0)>0\). In \eqref{eq:wave}, \(M=2\kappa Z-a''/a=-\gamma/t^2\) with
\[
 \gamma=p_0(p_0-1)-2\kappa\sigma/(1-q_0)^2,\qquad p_0=q_0/(1-q_0).
\]
Every real \(\gamma\) is realized by taking \(q_0\) sufficiently close to one and choosing
\(\kappa\sigma=[q_0(2q_0-1)-\gamma(1-q_0)^2]/2\), which tends to \(1/2\) while the upper bound tends to \(2\). The matter potentials may differ across models. Positive \(\gamma\) gives a negative lower-order canonical wave coefficient, not a kinetic ghost; no positive conserved energy of the evolving shear is asserted. This is a separate isotropic realization from the nonzero-Weyl Darboux examples, and no condition at the singular origin is used.

For a measurement interval in any finite positive-time slab, Theorem~\ref{thm:recover} recovers the load on its two-way shortened time square. The travel delays are in calibrated conformal time, not logarithmic time. The rank classification consequently applies to every separated pair inside that square. Earlier work on inverse-square couplings and integer Darboux potentials treats Huygens propagation in other space dimensions \cite{BL}. The claimed result here is the separated retarded load rank in the isolated plane channel, not Huygens propagation in the full spacetime.


\section{What the finite window decides}
A load recovered by Theorem~\ref{thm:recover} is a physical linear response in the specified channel. Its finite minors can exclude any proposed bound on real linear state dimension. The exact output on an open observation window depends only on the recovered load square, and on the known retained bulk. An exterior change with sufficient additional propagation delay is invisible. This restriction remains even for exact data.

For a sharp spectral example, take the calibrated constant positive potential \(m^2\). A half-line and a finite exterior of length \(\ell\) with a perfect Dirichlet reflector agree on all local area measurements whenever \(2\ell>T-2d\). Relative to a time window of length \(T\), the two loads differ only after a wall-to-reflector round trip. Their Laplace factors are respectively \(C p_m\) and \(C p_m\coth(\ell p_m)\), \(p_m=\sqrt{p^2+m^2}\) with positive real part on the appropriate half-plane. The half-line has continuous normal wavenumber; the finite interval has Dirichlet normal wavenumbers \(j\pi/\ell\). Thus exact finite-window retarded data can agree for an absolutely continuous and a discrete exterior force spectrum. The compact support of the auxiliary Fourier weight in \eqref{eq:besselFourier} is a property of an entire continuation of a tail; it is not the force spectrum.

For the nonstationary models we make no vacuum or stationary spectral assignment. A finite realization of a causal tail includes local derivative contact terms in the force law; it is not a construction of a closed lossless finite-dimensional Hamiltonian for the whole exterior or of its quantum state. Conversely, an infinite rank excludes every finite linear state realization under the stated comparison assumptions, but not approximation, nonlinear quantum descriptions, or a finite device with independent access to an uncounted propagating reservoir.

\begin{thebibliography}{15}
\raggedright
\bibitem{P44} AI Automated Quantum Gravity Research, \emph{Gravitational Region Sewing: Minimal Corner Registers and Continuous Radiative Loads}, Paper 44 (2026). \href{https://github.com/g3173028019-beep/ai-automated-quantum-gravity-research/tree/be3fc29a45b88b98fe5963888ed731a44d47ae59/papers/44-gravitational-region-sewing}{Source and PDF}.
\bibitem{P47} AI Automated Quantum Gravity Research, \emph{Spacelike area tomography and compact linear gravitational probes}, Paper 47 (2026). \href{https://github.com/g3173028019-beep/ai-automated-quantum-gravity-research/tree/be3fc29a45b88b98fe5963888ed731a44d47ae59/papers/47-spacelike-area-tomography}{Source and PDF}.
\bibitem{Khavkine} I. Khavkine, Local and gauge invariant observables in gravity, \emph{Class. Quantum Grav.} \textbf{32} (2015), 185019. \href{https://arxiv.org/abs/1503.03754}{arXiv:1503.03754}.
\bibitem{BLOS} L. Busch, M. Lassas, L. Oksanen and M. Salo, \emph{On Exponential Instability of an Inverse Problem for the Wave Equation}, \href{https://arxiv.org/abs/2506.23559}{arXiv:2506.23559}.
\bibitem{Kian} A. Feizmohammadi, J. Ilmavirta, Y. Kian and L. Oksanen, \emph{Recovery of time dependent coefficients from boundary data for hyperbolic equations}, \href{https://arxiv.org/abs/1901.04211}{arXiv:1901.04211}.
\bibitem{ENW} S. Endlich, A. Nicolis and J. Wang, Solid inflation, \emph{J. Cosmol. Astropart. Phys.} (2013) 10, 011. \href{https://arxiv.org/abs/1210.0569}{arXiv:1210.0569}.
\bibitem{KMT} V.~V. Kravchenko, S. Morelos and S.~M. Torba, Liouville transformation, analytic approximation of transmutation operators and solution of spectral problems, \emph{Appl. Math. Comput.} \textbf{273} (2016), 321--336. \href{https://arxiv.org/abs/1412.5237}{arXiv:1412.5237}.
\bibitem{Lions} J.-L. Lions, Supports de produits de composition. I, \emph{C. R. Acad. Sci. Paris} \textbf{232} (1951), 1530--1532; see also L. H\"ormander, Supports and singular supports of convolutions, \href{https://doi.org/10.1007/BF02391861}{\emph{Acta Math.} 110 (1963), 279--302}.
\bibitem{TVrank} A.-J. van der Veen and P.~M. Dewilde, On low-complexity approximation of matrices, \emph{Linear Algebra Appl.} \textbf{205--206} (1994), 1145--1201. Time-varying state-space rank.
\bibitem{PT} P. Gaillard and V.~B. Matveev, Wronskian addition formula and Darboux--P\"oschl--Teller potentials, \href{https://doi.org/10.1155/2013/645752}{Journal of Mathematics (2013), 645752}.
\bibitem{Legendre} P.~J. Zeitsch, On the Riemann function, \emph{Mathematics} \textbf{6} (2018), 316, Eq.~(24). \href{https://doi.org/10.3390/math6120316}{doi:10.3390/math6120316}.
\bibitem{BL} Y. Berest and I. Loutsenko, Huygens' principle in Minkowski spaces and soliton solutions of the Korteweg--de Vries equation, \emph{Comm. Math. Phys.} \textbf{190} (1997), 113--132. \href{https://arxiv.org/abs/solv-int/9704012}{arXiv:solv-int/9704012}.
\end{thebibliography}
\end{document}
